\section{The mechanical question and its interfaces}
Let $a=(1,0)$, $b=(1/2,\sqrt3/2)$, $\Lambda=\Z a+\Z b$, and
\[
 I=[9/20,47/100],\qquad Q_R=(\R^2/\Lambda)\setminus\overline B(0,R).
\]
The speed is one and every reflection is specular. Three indices will be kept distinct: a physical collision count $N$, an induced return count $n$, and the number of experimental bits. An induced record has values in $\Z^2\times\Z\times\R$, with counting measure in the first three coordinates and Lebesgue measure in the last. A four-dimensional Lebesgue density would be the wrong object.

The geometric inverse of the preceding article identifies the table from two fixed collision fields \cite{A2G}. It supplies neither independent flights nor an induced transfer-operator estimate. Conversely, the stationary law of the equilibrium billiard is not the unknown preparation law of that inverse experiment. The distinction permits geometric estimates to be used for equilibrium statistics without first estimating the entire preparation law.

The joint arithmetic result is Theorem~\ref{thm:joint-periodic}. It concerns
the complete periodic phase of an explicitly enlarged physical first-return
section. Lemma~\ref{lem:excursion-family} constructs its infinite period
family, and Theorem~\ref{thm:period-excess} proves the strict excess
inequalities. These results go beyond the zero-roof determinant by ruling
out every nonzero roof frequency, without a numerical irrationality claim.
The real-rank and compact-frequency consequences are stated separately
from measurable cohomology regularity and operator contraction.

The original physical record, winding geometry, critical-edge and
conditional inversion arguments remain in the article. Theorem
\ref{thm:local-edge-LLT} refines the inversion requirement by using local
edge errors instead of small total edge mass. Propositions
\ref{prop:exact-induced-means} and \ref{prop:induced-suspension} give exact
normalizations and stationary endpoint control for the enlarged section.
The full raw local limit is not inferred from the arithmetic theorem:
its spectral, residual and weighted-insertion requirements are listed in
the final section. All probability laws here refer to actual specular
orbits, not independent flight surrogates.

The relevant prior local-limit theory includes the Lorentz-process theorem of Sz\'asz--Varj\'u \cite{SV} and the suspension framework of Dolgopyat--N\'andori \cite{DN}. Definition 2.2 and Proposition 6.3 of the latter concern a mixing local limit tested against fixed compactly supported functions under minimality and group assumptions. They do not state absolute integrability of every raw finite-time characteristic function, or the particular geometrically uniform density theorem sought here. Demers--Zhang \cite{DZ} give a perturbation framework for Lorentz-gas maps; its spectral continuity is not automatically twice differentiability in a moving geometric parameter. We use the corrected version of the flow-mixing work of Baladi--Demers--Liverani \cite{BDL}. No novelty is claimed for Fourier inversion, the Morse lemma, or the suspension identity as general methods.


Theorem~\ref{thm:return-stability} compares the actual moving return
events in density $L^1$ and first moment, at every fixed return count.
Corollary~\ref{cor:uniform-stationary-endpoints} then supplies uniform
stationary endpoint control over the whole radius interval. The proof uses
finite-itinerary stability, coarea and exact Kac means, not a hypothetical
common spectral space. These finite-record moduli do not purport to supply
the quantitative long-time estimates in the local-limit interface.

Theorem~\ref{thm:quantitative-excess} strengthens strict increase to
uniform two-sided exponential bounds, obtained from the exact nonlinear
half action. Theorem~\ref{thm:log-period-separation} uses their upper and
lower bounds together: the former selects periods of logarithmic length,
and the latter supplies a polynomial discrepancy rather than mere
nonvanishing. Corollary~\ref{cor:approximate-phase-bound} states precisely
the continuous-phase consequence. It does not assume the measurable
regularity needed for an operator application.

At the level of the clock, Theorems~\ref{thm:uniform-return-fluid} and
\ref{thm:uniform-physical-clock} turn fixed-record continuity and exact
means into a uniform functional first-order law. This includes the maximum
unfinished return over the observation window at scale $t$, whereas the
retained stationary endpoint result is at scale $\sqrt t$ for finitely
many deterministic endpoints. The distinction of scales is essential.
Corollary~\ref{cor:physical-rate-error} gives a geometric error term for
the physical rate without assuming independence of the radius estimate
and the observed orbit. No originality is claimed for the general ergodic,
subadditive or renewal inversion arguments; the point is their verified
application to the actual moving return family.
